\documentclass[11pt]{amsart}
\usepackage{amsmath,amssymb,amsthm}
\usepackage[margin=1.15in]{geometry}
\usepackage{hyperref}
\usepackage{booktabs}

\theoremstyle{plain}
\newtheorem{theorem}{Theorem}[section]
\newtheorem{lemma}[theorem]{Lemma}
\newtheorem{proposition}[theorem]{Proposition}
\newtheorem{corollary}[theorem]{Corollary}
\theoremstyle{definition}
\newtheorem{definition}[theorem]{Definition}
\newtheorem{remark}[theorem]{Remark}
\newtheorem{observation}[theorem]{Observation}
\newtheorem{conjecture}[theorem]{Conjecture}

\newcommand{\Norm}[1]{\left\lVert #1 \right\rVert}
\newcommand{\ML}{\operatorname{ML}}
\newcommand{\cpt}{\operatorname{cap}}
\newcommand{\srp}{\operatorname{srp}}

\begin{document}

\title[A verified covering framework for the lonely runner problem]
{A verified covering framework for the lonely runner problem in bounded windows}

\author{Vico Bonfioli}
\email{vicobonfioli@gmail.com}

\date{18 October 2026. Preprint, not peer reviewed.}

\begin{abstract}
We develop, and verify in Lean 4, a covering framework for the lonely runner conjecture restricted
to bounded velocity windows. Everything labelled verified below is formalized with no axioms beyond
\texttt{propext}, \texttt{Classical.choice} and \texttt{Quot.sound}. First, the window statement is
equivalent to the assertion that a hitting number attains a Dirichlet ceiling; this recasting is
known in substance and we prove both directions inside the proof assistant. Second, we introduce a
surplus normalization of the cover-capacity certificate whose budget is the constant $2$,
independent of the number of velocities, and prove it equivalent to the usual capacity condition.
Third, we classify the bad sets of size at most three into two singleton families, four pair
families and three triple families, and prove each of them; the singleton and pair families are
exhaustive for every $m \le 260$, the triple families are not, and we exhibit minimal triples
outside them. We also prove that the complementary pair family is confined to exactly two moduli.
Fourth, we prove a dichotomy on residues: if two velocities are simultaneously close to an integer
multiple of the period, then so is their difference or their sum. It yields two of the three shapes
a bad pair can take, uniformly in $m$, and the third follows from an identity between the two
residues which we prove by a lattice argument, so that the classification is complete: a bad pair is
$\{a,2a\}$ or has $a+b \in \{2m+1,2m+2\}$ with $\gcd(a,a+b)=1$. It also shows that the top
third of the window carries no bad set, so the bad-pair graph has an independent set of size $m-1$
for every $m$.

As an application of the associated decision procedure we decide the window $\{1,\dots,2m\}$ by
exhaustive computation for every $m \le 63$, extending to that range the statement previously
verified for $m \le 24$; the extension is computational and is not itself formalized. The same range
is then obtained a second time from the nine families alone, with the exhaustive enumeration
discarded, so that every set entering the search is one that we have proved to be a bad set. We
emphasize that all results are restricted to velocities bounded by $2m$; the unrestricted conjecture
is known for at most thirteen runners by different methods, and the two are not comparable. We also
record five closed approaches to the uniform statement, including a counterexample at $m = 19$ to
the block template suggested by the smaller cases, and counterexamples at $m = 150, 318, 486, 570$
to a regularity of the bad-pair graph that holds at every one of the thirty-three composite moduli
with $m \le 63$.
\end{abstract}

Statements are tagged with the status of their proof: \emph{Verified} means formalized in Lean~4
against Mathlib with no \texttt{sorry} and no axioms beyond \texttt{propext},
\texttt{Classical.choice} and \texttt{Quot.sound}; \emph{Computed} means established by a finite
computation, with the generating script cited. An untagged theorem, lemma, proposition or corollary
carries a complete proof in the text and no formalization. Conjectures are labelled as such and are
proved nowhere here. Definitions carry no status.


\maketitle

\section{Introduction}

Let $v_1,\dots,v_m$ be distinct positive integers and write $\Norm{x}$ for the distance from $x$ to
the nearest integer. Put
\[
  \ML(v_1,\dots,v_m) \;=\; \sup_{t \in \mathbb{R}} \ \min_{1 \le i \le m} \Norm{v_i t}.
\]
The lonely runner conjecture, due to Wills and independently to Cusick in the language of view
obstruction, asserts
\begin{equation}\label{eq:lrc}
  \ML(v_1,\dots,v_m) \;\ge\; \frac{1}{m+1}
\end{equation}
for every such tuple. In the runner formulation there are $m+1$ runners, $m$ velocities relative to
the observer. The bound is sharp: $v_i = i$ gives equality, and Dirichlet's approximation theorem
shows $\ML(1,\dots,m) \le 1/(m+1)$, so \eqref{eq:lrc} is the statement that a Dirichlet ceiling is
attained. See Perarnau and Serra \cite{PS} for a survey.

The conjecture is settled for small numbers of runners. At the time of writing it is known for at
most thirteen runners, equivalently for $m \le 12$ velocities in the normalization above, by
Rosenfeld \cite{Ros8,Ros9}, Trakulthongchai \cite{Tra}, and Sungkawichai and Trakulthongchai
\cite{ST}; all of these are computer assisted and descend from a prime divisibility criterion. We
state the convention explicitly because the literature is divided between counting runners and
counting velocities.

A separate line of work restricts the velocities to a window. Pomerance \cite{Pom} states that it
remains a challenge to prove \eqref{eq:lrc} when $\{v_1,\dots,v_m\} \subseteq \{1,\dots,2m\}$, and
the same restricted problem appears as Problem 6 of \cite{PS}. All results in this paper concern
that restricted problem, and we state the restriction with every claim. The restricted and
unrestricted problems are not comparable, and the method here is not competitive on the
unrestricted one; see Remark~\ref{rem:regime}.

Substantial parts of this window are already closed, and we state which, since our contribution is
what remains. Tao \cite{Tao} proved that $v_m < 1.2m$ implies
\eqref{eq:lrc}, uniformly in $m$ \cite[Theorem 22]{PS}; this closes the lower three fifths of the window outright and is
the result that prompts the question of raising $1.2$ to $2$. Bohman and Peng \cite{BP}, via coprime
mappings, proved that for sufficiently large $m$ the range $m < v_m \le 2m - \exp(c(\log\log m)^2)$
is also closed, and Pomerance \cite{Pom} improved the error term to $c(\log m)^2$; both are
ineffective in $m$ and so say nothing about any concrete $m$. Pandey \cite{Pan} proved
\eqref{eq:lrc} for every $m$-subset of an arithmetic progression $\{a, a+d, \dots, a+(k-1)d\}$ of
length $k \le 2m-3$ with $a$ and $d$ not both $1$; taking $d = 1$ this closes every velocity set
inside the window whose least element exceeds $1$ and whose span $v_m - v_1 + 1$ is at most $2m-3$.

What survives all of these at a given $m$ is the set of $V \subseteq [1,2m]$ with $v_m \ge 1.2m$ and
with either $v_1 = 1$ or $v_m - v_1 + 1 \ge 2m-2$. That residual contains the extremal
configurations: the tight instance $\{1,\dots,m\}$ has $v_1 = 1$, and the tight instances with
$v_m = 2m - \Theta(\log m)$ noted in \cite[\S8.2]{PS} nearly span the window. The results above are
all of the form ``\eqref{eq:lrc} holds for velocity sets satisfying a structural condition''. None
of them decides the window at a fixed $m$, that is, settles \eqref{eq:lrc} for \emph{every}
$m$-subset of $[1,2m]$ simultaneously, and it is that decision which the present method performs.

\subsection{Results}

Throughout, a claim is labelled \emph{verified} if it is formalized in Lean 4 with a clean build,
no \texttt{sorry}, and no axioms beyond \texttt{propext}, \texttt{Classical.choice} and
\texttt{Quot.sound}; \emph{proved} if a complete human proof is given but no formalization;
\emph{computed} if it rests on an exhaustive machine enumeration that is not itself formalized;
and \emph{conjectural} otherwise. Section~\ref{sec:lean} gives the correspondence between numbered
statements and Lean declarations.

Fix $m \ge 1$ and let $t = m+1$. For coprime integers $p,q$ with $q > 0$ define the \emph{bad set}
of the time $p/q$ inside a window $[1,N]$ by
\[
  B_N(p,q) \;=\; \bigl\{\, x \in [1,N] \ :\ \Norm{x p / q} < 1/t \,\bigr\},
\]
and let $\tau(N)$ be the least cardinality of a subset of $[1,N]$ meeting every $B_N(p,q)$.

\begin{theorem}[Verified; see Section~\ref{sec:hitting}]\label{thm:main-equiv}
For $N \ge m+1$, inequality \eqref{eq:lrc} holds for every $m$-element velocity set contained in
$[1,N]$ if and only if $\tau(N) > m$. Moreover $\tau(N) \le m+1$ for every $N$, so the window
statement is equivalent to $\tau(N) = m+1$ exactly.
\end{theorem}

The content of Theorem~\ref{thm:main-equiv} is known. The covering formulation of \eqref{eq:lrc} is
standard \cite{Tao}; that an optimal time may be taken of the form $\ell/(v_i+v_j)$ is elementary
and is used throughout the literature without attribution \cite[eq.~(4)]{PS}, while the resulting
replacement of the supremum by a finite maximum is credited to Haralambis \cite[Remark 1]{Har},
with a proof in \cite[Theorem 6]{CG}; and Rosenfeld \cite[\S6, Lemma 7]{Ros8} already runs exactly
the resulting finite set-cover search inside a velocity window. The upper bound $\tau(N) \le m+1$ is
the classical Dirichlet observation recorded in \cite[\S4]{PS}. Our contribution to
Theorem~\ref{thm:main-equiv} is the formalization and the naming of $\tau$, not the mathematics.

The next statement is a normalization of the cover-capacity certificate of
Section~\ref{sec:surplus}. It is a change of variables and we claim nothing for it beyond the
bookkeeping: the point is that the threshold becomes the constant $2$ instead of the $m$-dependent
$m-1$, which is what makes the parts reusable across $m$.

\begin{theorem}[Verified; Lemma~\ref{thm:surplus}]\label{thm:main-surplus}
Let $U_1,\dots,U_r$ cover $[1,2m]$, and suppose that for each $i$ every subset of $U_i$ with more
than $c_i$ elements contains a bad set. If $\sum_i |U_i| = 2m$ and
\[
  \sum_{i=1}^{r} \srp_i \;\ge\; 2, \qquad \srp_i := |U_i| - 2c_i,
\]
then \eqref{eq:lrc} holds for every $m$-element velocity set contained in $[1,2m]$.
\end{theorem}

The point is the quantifier. The usual condition $\sum_i c_i < m$ has a budget depending on $m$, so
a certificate must be rebuilt for each $m$; the surplus budget is the constant $2$ at every $m$. The
two singleton bad sets supply exactly two units of surplus, and by Proposition~\ref{prop:surplus-one}
they are the entire budget available.

\begin{theorem}[Computed; Families verified, exhaustiveness computed; Section~\ref{sec:vocab}]\label{thm:main-vocab}
Nine explicitly parameterized families of bad sets are proved: two of singletons, four of pairs and
three of triples. The complementary pair family occurs at exactly two moduli, namely $q = 2m+1$ and
$q = 2m+2$, and at no others. The pair families are exhaustive for \emph{every} $m$: an exact bad
pair $a<b$ satisfies $b = 2a$ or $a+b > 2m$, so it is $\mathrm{DBL}$ or $\mathrm{CMP}$ (Lean:
\texttt{badpair\_doubling\_or\_large}, with no hypothesis on the modulus). For singletons the
exhaustiveness is computed, over $4 \le m \le 260$. The triple families are
\emph{not} exhaustive, and Proposition~\ref{prop:triples-incomplete} exhibits minimal bad triples
outside them.
\end{theorem}

The families themselves, the two-moduli confinement and the exhaustiveness at size two are
verified; the exhaustiveness at size one is computed, over a range wider than any other computation
in this paper. The size-two half is unconditional in the modulus, which is why it needs no range at
all; the size-one half still rests on a sweep. We separate the two sizes because the distinction is
load-bearing: the argument that a particular element closes
the independence gap uses only the pair families, whereas the decision procedure of
Theorem~\ref{thm:fam63} uses all nine and, as explained in Section~\ref{sec:vocab}, is unaffected by
the triples that the classification misses.

\begin{theorem}[Computed; Section~\ref{sec:computation}]\label{thm:main-63}
For every $m$ with $4 \le m \le 63$, inequality \eqref{eq:lrc} holds for every $m$-element velocity
set contained in $[1,2m]$.
\end{theorem}

The previously verified range for this restricted problem was $m \le 24$. Theorem~\ref{thm:main-63}
is obtained from Theorem~\ref{thm:main-equiv} by deciding $\tau(2m) > m$ directly, which is a
decision problem rather than a search for a certificate, and is correspondingly cheaper. The
verdicts for $m \le 24$ reproduce the previously verified range, and the range $25 \le m \le 63$ is
new. We stress that Theorem~\ref{thm:main-63} is computed and not verified: lifting it to a Lean
proof requires a certificate for each $m$, and producing one is an open problem for each $m$
(Remark~\ref{rem:lift}).

The computation admits a sharper description. Because $\tau$ is monotone in the family of sets to be
hit, the decision may be run against the nine families of Theorem~\ref{thm:main-vocab} alone, with
the exhaustive enumeration discarded; it then still exhausts for every $m \le 63$
(Theorem~\ref{thm:fam63}). On that route the exhaustiveness clause of Theorem~\ref{thm:main-vocab}
is not used, every set fed to the search is one of the verified families, and the only unformalized
ingredient left in Theorem~\ref{thm:main-63} is the finite search itself. That every generated set
really is a bad set, which is what soundness needs, is certified constructively for all
$m \le 250$; see Section~\ref{sec:computation}.

Finally, Section~\ref{sec:obstructions} records five approaches to the uniform statement that are
now closed, with the reasons. Among them is a counterexample at $m = 19$ to the block template that
the cases $m = 10,13,16$ suggest.

\subsection{What is not claimed}

The uniform statement, that \eqref{eq:lrc} holds for every $m$ and every velocity set inside
$[1,2m]$, remains open and nothing here approaches it. The framework certifies individual $m$; it
has no known uniform selection mechanism, and Section~\ref{sec:obstructions} explains why the
natural candidates fail.

\begin{remark}\label{rem:regime}
The efficiency of the method is governed by the minimum bad-set size, which by
Theorem~\ref{thm:mincount} equals $\lfloor N/(m+1) \rfloor$. In a window $N = Bm$ with $B$ bounded
this is a small constant and the decision procedure is fast; when $N/m$ grows the minimum bad set
occupies a constant fraction of the window and the procedure degenerates. Applying the window
reduction of \cite{MSS}, which permits taking $N = \binom{m+1}{2}^{m-1}$, the method reproves the
unrestricted conjecture for at most four runners and is already impractical at five. It is
therefore far behind the prime divisibility sieve on the unrestricted problem, and its regime is
the bounded window.
\end{remark}

\section{Setup}\label{sec:setup}

We work with the integer form of the badness predicate, which is decidable and is what the
formalization manipulates.

\begin{definition}\label{def:good}
For $k, x, p, q \in \mathbb{N}$ with $q > 0$ set
\[
  \operatorname{good}_k(x,p,q) \;=\;
  \bigl[\, q \le (k+2)\min\bigl(xp \bmod q,\ q - (xp \bmod q)\bigr) \,\bigr].
\]
We say $x$ is \emph{bad} at $(p,q)$ when $\operatorname{good}_k(x,p,q)$ fails.
\end{definition}

Writing $m = k+1$, badness at $(p,q)$ is exactly $\Norm{xp/q} < 1/(m+1)$, so
$B_N(p,q) = \{x \in [1,N] : x \text{ bad at } (p,q)\}$. A velocity set $V$ with $|V| = m$ fails
\eqref{eq:lrc} precisely when every bad set meets $V$, because the supremum defining $\ML$ is
attained at a rational time whose denominator is a pairwise sum of velocities
\cite[eq.~(4)]{PS}, \cite{Har,CG}. This is the input to
Theorem~\ref{thm:main-equiv}.

The following counting bound is used throughout and is sharp.

\begin{theorem}[Verified]\label{thm:mincount}
For all $k, N$ and all $p,q$ with $q > 0$,
\[
  \bigl|\{x \in [1,N] : x \text{ bad at } (p,q)\}\bigr| \;\ge\; \Bigl\lfloor \frac{N}{k+2} \Bigr\rfloor .
\]
\end{theorem}

\begin{proof}
Distribute the $N+1$ points $\{0 \cdot p/q\}, \dots, \{N \cdot p/q\}$ among the $k+2$ half-open
boxes $[j/(k+2), (j+1)/(k+2))$. Some box contains at least $\lfloor N/(k+2)\rfloor + 1$ of them.
Subtracting the least index in that box from the others gives $\lfloor N/(k+2)\rfloor$ distinct
values in $[1,N]$, each bad, since two points of a common half-open box of length $1/(k+2)$ differ
by less than $1/(k+2)$.
\end{proof}

Theorem~\ref{thm:mincount} is classical. It is the Dirichlet box principle with $B$ points per box
in place of two, and also follows from van der Corput's generalization of Minkowski's theorem
applied to $\{|x| < N + \delta,\ |y| < 1/(k+2)\}$ and the unimodular lattice $\{(v, vp/q - u)\}$;
see \cite{Cassels}. The same inequality appears inside a proof in \cite[eq.~(1.34)]{BHV}, there
under a linear independence hypothesis that excludes the rational times at which it is sharp. It is
sharp: at $p/q = 1/(k+2)$ the bad set is the set of multiples of $k+2$.

\section{The hitting formulation}\label{sec:hitting}

\begin{theorem}[Verified]\label{thm:ceiling}
Every bad set meets $[1,k+2]$. Consequently $\tau(N) \le k+2 = m+1$ for every $N \ge m+1$.
\end{theorem}

\begin{proof}
Apply Theorem~\ref{thm:mincount} with $N = k+2$; the floor is $1$, so the bad set of any time has an
element in $[1,k+2]$. Hence $[1,m+1]$ is itself a hitting set.
\end{proof}

\begin{theorem}[Verified]\label{thm:equiv}
Fix $N \ge m+1$.
\begin{enumerate}
\item If every subset of $[1,N]$ with at most $m$ elements misses some bad set, then every
  $m$-element velocity set inside $[1,N]$ satisfies \eqref{eq:lrc}.
\item Conversely, if some $T \subseteq [1,N]$ with $|T| \le m$ meets every bad set, then there are
  $m$ distinct positive velocities inside $[1,N]$ violating \eqref{eq:lrc}.
\end{enumerate}
In particular the window statement for $[1,N]$ is equivalent to $\tau(N) > m$, and with
Theorem~\ref{thm:ceiling} to $\tau(N) = m+1$.
\end{theorem}

\begin{proof}
For (1): a velocity set violating \eqref{eq:lrc} meets every bad set, so it is itself a hitting set
of size $m$; the hypothesis applied to it produces a time at which every velocity is good, which
contradicts the violation. For (2): extend $T$ to exactly $m$ distinct elements of $[1,N]$, which
is possible since $N \ge m$; the extension still meets every bad set, and a family meeting every bad
set has no good time, hence violates \eqref{eq:lrc}, by the attainment of the supremum defining
$\ML$ at a rational time with denominator a pairwise sum of velocities, recorded above with
references \cite[eq.~(4)]{PS}, \cite{Har,CG}.
\end{proof}

\begin{remark}
The margin in Theorem~\ref{thm:equiv} is always exactly one, by Theorem~\ref{thm:ceiling}. This is
the tight instance $\{1,\dots,m\}$ restated in hitting language and should not be read as a
structural discovery; the ceiling is classical \cite[\S4]{PS}. It does explain why every certificate
in this framework is tight, which is otherwise a recurring surprise.
\end{remark}

\section{The surplus certificate}\label{sec:surplus}

For $U \subseteq [1,2m]$ let $\cpt(U)$ denote the largest size of a subset of $U$ containing no bad
set, equivalently $|U| - \tau(U)$ where $\tau(U)$ is the least size of a subset of $U$ meeting every
bad set contained in $U$.

The cover-capacity certificate, in the form we use, states that if $U_1,\dots,U_r$ cover $[1,2m]$
and every subset of $U_i$ with more than $c_i$ elements contains a bad set, then $\sum_i c_i < m$
implies the window statement. The mechanism is that the complement of a violating velocity set is
an independent set of size $m$ inside the window, and the parts bound it piecewise. At a general
window height $Bm$ the complement has $(B-1)m$ elements and the condition becomes
$\sum_i c_i < (B-1)m$; both forms are verified. The soundness of testing a weighted combination of times, equivalently a finitely supported measure, against a family of velocities is Bedert's \cite[Lemma~2.1]{Bed}; we claim no priority for it.

Note that the $c_i$ are upper bounds supplied with the parts, so $c_i \ge \cpt(U_i)$, with equality
only when the bound is exact; the surplus below is computed from the $c_i$ and not from $\cpt$.
Note also that $\sum_i |U_i| = 2m$ together with the covering condition forces the $U_i$ to
partition $[1,2m]$.

\begin{lemma}[Verified]\label{thm:surplus}
Set $\srp_i := |U_i| - 2c_i$. Under the hypotheses above with $\sum_i |U_i| = 2m$, the condition
$\sum_i c_i < m$ is equivalent to $\sum_i \srp_i \ge 2$.
\end{lemma}

\begin{proof}
$\sum_i \srp_i = \sum_i |U_i| - 2\sum_i c_i = 2m - 2\sum_i c_i$, and $2m - 2\sum_i c_i \ge 2$ is
$\sum_i c_i \le m-1$.
\end{proof}

This is arithmetic, not mathematics: it is the same condition written so that the threshold no
longer depends on $m$. We state it separately only because the search consumes it in that form.

The identity is immediate; its value is that the budget no longer depends on $m$. A certificate
becomes: cover the window by parts of nonnegative surplus and locate two units of surplus.

\begin{proposition}[Proved]\label{prop:surplus-one}
For every $m\ge1$ the only bad sets of size one inside $[1,2m]$ are $\{1\}$ and $\{m+1\}$.
\end{proposition}

\begin{proof}
Fix $\alpha$ and let $\varepsilon=1/(m+1)$. Order the $2m+1$ points $\{x\alpha\}$, $x=0,\dots,2m$,
on the circle and let $s$ be the least index difference realised by an adjacent pair, with gap $u$.
The gaps sum to $1$; every gap whose index difference is not $s$ has length at least $\varepsilon$,
since a shorter one would exhibit a bad singleton of smaller index difference; and at most
$2m+1-s$ gaps have difference $s$. Hence $s\varepsilon\le1$, so $s\le m+1$.

If $s\le m$ then $2s$ lies in the window, and $2s$ good forces $u\ge\varepsilon/2$. The points are
then distinct and every pair $\{x,x+s\}$ is adjacent with gap exactly $u$, so
$1\ge(2m+1-s)\varepsilon/2+s\varepsilon$, which forces $s=1$. Thus $s\in\{1,m+1\}$, and
$s=m+1$ requires $(m+1)\alpha\in\mathbb{Z}$. Both are realised, at $\alpha$ near $0$ and at
$\alpha=1/(m+1)$ respectively. No approximation theorem is used; the count is the mechanism behind
the three-distance theorem. Verified independently by exhaustive breakpoint sweep in exact integer
arithmetic for $2\le m\le400$, with a negative control at threshold $1/(m+2)$ that fails at every
$m$.
\end{proof}

\begin{proposition}[Computed]\label{prop:surplus-two}
The number of bad sets of the minimum size $B-1$ inside $[1,Bm]$ is exactly two for every
$B \in \{3,4,5\}$ and every $m$ in the computed range. (The range of $m$ is not recorded and the generating program is not archived; HEURISTIC.) Together with
Proposition~\ref{prop:surplus-one}, which settles $B=2$ for every $m$, the surplus available from
minimum-size parts is exactly two.
\end{proposition}

Parts of nonnegative surplus known to us are: a singleton bad set, surplus $+1$; a bad pair,
surplus $0$; an odd cycle of bad pairs of length $2\ell+1$, whose vertex cover number is $\ell+1$
and independence number $\ell$, surplus $+1$; and a four-element part all four of whose triples are
bad, surplus $0$.

\begin{proposition}[Verified]\label{prop:packing}
Let $S_1,\dots,S_{m+1}$ be pairwise disjoint subsets of $[1,Bm]$, each of size at least $B-1$. Then
at least $B$ of them have size exactly $B-1$.
\end{proposition}

\begin{proof}
Let $a$ be the number of parts of size exactly $B-1$ and $b$ the number of larger parts, so
$a + b = m+1$ and $(B-1)a + Bb \le Bm$. Substituting $a = m+1-b$ gives $b \le m-B+1$, hence
$a \ge B$.
\end{proof}

Proposition~\ref{prop:packing} combined with Theorem~\ref{thm:mincount} and
Proposition~\ref{prop:surplus-two} shows that a certificate assembled from pairwise disjoint bad
sets exists only at $B = 2$: it would require $B$ minimum-size bad sets, and only two exist. At
$B = 2$ the requirement is exactly two and exactly two are available, which is why the dense-window
certificates are tight and why the two free singletons are structurally necessary rather than
convenient.

\section{The vocabulary of small bad sets}\label{sec:vocab}

Fix $m$ and write $k = m-1$. We classify the bad sets of size at most three inside $[1,2m]$. A bad
set is \emph{minimal} if it contains no bad set properly.

The enumeration is exhaustive over all real times, by the following observation. Membership of a
fixed $x$ in the bad set changes only where $\Norm{x\alpha} = 1/(m+1)$, that is at
$\alpha = (j(m+1) \pm 1)/(x(m+1))$, a rational of denominator dividing $x(m+1) \le N(m+1)$. Between
two consecutive such breakpoints $a/b < c/d$ in lowest terms the mediant $(a+c)/(b+d)$ lies strictly
between and has denominator at most $2N(m+1)$. Hence scanning every reduced $p/q$ with
$q \le 2N(m+1)$ meets every breakpoint and every interval between consecutive breakpoints, and the
resulting list of bad sets is the list occurring for any real time.

\begin{theorem}[Computed]\label{thm:oth}
For every $m$ with $4 \le m \le 260$, every minimal bad set of size one or two inside $[1,2m]$
belongs to one of the families $\{1\}$, $\{m+1\}$, MUL, CMP, DBL below. Over that range this was
checked against every bad set of size at most three produced by an enumeration exhaustive over all
real times, and no further species of size one or two occurs. The program is \texttt{code/covering/clsweep.rs}: it checks $2\,341\,337$ minimal bad sets of size at most three over $4\le m\le260$ (the sets of size three that lie in none of the families are the triples of Proposition~\ref{prop:triples-incomplete}) and reports no unclassified singleton or pair. An earlier pair-only program, \texttt{code/covering/paircls.rs}, gives the same verdict for $4\le m\le32$; it stores a pair in a $64$-bit mask and is void for $m\ge33$.
\end{theorem}

For triples the corresponding statement is \emph{false}, and we record it as such.

\begin{proposition}[Computed]\label{prop:triples-incomplete}
The families SUM, SUM$'$, DBL$_3$ do not exhaust the minimal bad triples. Roughly one minimal bad
triple in eight lies outside them, at every $m$ from $7$ to $260$. Two smallest examples:
$\{5,8,13\}$ at $m = 9$, realized at $61/160$, and $\{6,12,18\}$ at $m = 10$, realized at
$34/209$; both are minimal, and neither satisfies the hypotheses of any of the three propositions.
The program \texttt{code/covering/bad\_triple\_classify.rs} enumerates all bad sets of size at most three over all real times and compares them as sets with the predicted list, with $0$ discrepancies, for $4\le m\le260$: it finds $2\,272\,034$ minimal triples, of which $294\,170$ lie outside the three families (types S3 and CHN of the exact classification), at least one outside at every $m$ from $7$ to $260$, with a fraction between $0.12$ and $0.19$ at each such $m$.
\end{proposition}

The two failures are of different kinds and both are instructive. The triple $\{5,8,13\}$ has the
shape $\{a,b,a+b\}$ with $a,b$ coprime, but sits on the boundary of family SUM, where
$2a+b = 18$ is not strictly greater than $2m = 18$, and it is too low for SUM$'$, which needs
$b \ge m+1$. The triple $\{6,12,18\}$ is the set of multiples of $6$ inside the window, and it has
both the SUM shape and the DBL$_3$ shape, but every family that could carry it requires a modular
inverse and $\gcd(6,12) = 6$. A classification by shape alone would count both as covered; the
families are defined by their hypotheses, and it is the hypotheses that fail.

We can prove two of the three shapes for pairs, uniformly in $m$, and it is worth isolating what
the obstruction to the third one is.

\begin{theorem}[Verified]\label{thm:pairdich}
Let $a < b$ both be bad at the time $p/q$. Then $b-a$ is bad or $a+b$ is bad.
\end{theorem}

\begin{proof}
Write $r_x = xp \bmod q$ and regard $r_x$ as a point on a circle of circumference $q$; badness of
$x$ says $r_x$ lies within $q/(m+1)$ of $0$, on one side or the other. If $r_a$ and $r_b$ lie on the
same side, the distance from $0$ of $r_b - r_a$ is at most the larger of the two distances, so
$b-a$ is bad; if on opposite sides, the same holds for $r_a + r_b$, so $a+b$ is bad.
\end{proof}

\begin{corollary}[Verified]\label{cor:pairstruct}
If the bad set of $p/q$ inside $[1,2m]$ is exactly $\{a,b\}$ with $1 \le a < b \le 2m$, then
$b = 2a$ or $a + b > 2m$.
\end{corollary}

\begin{proof}
By Theorem~\ref{thm:pairdich} one of $b-a$, $a+b$ is bad. In the first case $b-a$ lies in $[1,2m]$
and is smaller than $b$, so it equals $a$ and $b = 2a$. In the second, $a+b$ exceeds $b$, so it
cannot be one of the two bad values, and therefore cannot lie in the window at all.
\end{proof}

This is the doubling shape and the lower half of the complementary shape. The remaining gap is the
upper bound $a+b \le 2m+2$ \emph{in the non-doubling case}, and the restriction is essential rather
than cosmetic: $\{m,2m\}$ is a bad set for every $m$, by Proposition~\ref{prop:mul} at $d = m$, and
it has $a+b = 3m$, which exceeds $2m+2$ as soon as $m \ge 3$. Doubling pairs are already covered by
the first disjunct and must be excluded from the second.

That remaining bound holds as soon as the pair is realized at one particular modulus.

\begin{theorem}[Verified]\label{thm:pairupper}
Let $a < b \le 2m$ with $b \ne 2a$, let $ap \equiv 1 \pmod{a+b}$, and suppose the bad set of
$p/(a+b)$ inside $[1,2m]$ is exactly $\{a,b\}$. Then $a + b \le 2m+2$.
\end{theorem}

\begin{proof}
Put $q = a+b$. Since $b = q - a$, its residue is $-1$, and $b - a = q - 2a$ has residue $-2$. So
$b-a$ is bad exactly when $2(m+1) < q$. But $1 \le b - a \le 2m$, and $b - a$ is neither $a$, which
would force $b = 2a$, nor $b$, which would force $a = 0$; so $b-a$ bad contradicts the bad set being
$\{a,b\}$. Hence $q \le 2m+2$.
\end{proof}

The hypothesis that the pair is realized at that particular modulus can be weakened to an identity
at whatever modulus is given. Fix a time $p/q$ whose bad set inside $[1,2m]$ is exactly $\{a,b\}$,
and suppose the two residues lie on opposite sides of zero: write $u$ for the residue of $ap$ modulo
$q$ and $q - v$ for the residue of $bp$, with $0 < u$ and $0 < v < q$.

\begin{theorem}[Verified]\label{thm:pairident}
With that notation, if
\[
  av + bu = q,
\]
then $a + b \le 2m+2$.
\end{theorem}

\begin{proof}
Write $t = m+1$. Since $b \ne 2a$ and $a \ge 1$, the element $b-a$ lies in $[1,2m]$ and equals
neither $a$ nor $b$, so it is good. Its residue is $-(u+v)$, whence $q \le (u+v)t$. The identity
then gives $a(u+v) < av + bu = q \le (u+v)t$, so $a \le t - 1 = m$ and therefore $2a \le 2m$. As
$2a$ is neither $a$ nor $b$, it is good as well, and its residue is $2u$, whence $q \le 2ut$.
Writing $b = a+d$ with $d \ge 1$,
\[
  q(2a+d) = 2aq + dq \le 2a(u+v)t + d\cdot 2ut = 2t\bigl(av + (a+d)u\bigr) = 2tq ,
\]
and cancelling $q$ gives $a + b = 2a + d \le 2t = 2m+2$.
\end{proof}

Neither goodness fact is an extra hypothesis; both are forced by $\{a,b\}$ being the whole bad set,
and both are derived from that in the formalization. Nothing is required of $v$ beyond positivity.
The symmetric partner $q \le 2vt$ of the second fact is false, failing in $3280$ of the $11528$
instances below (program \texttt{code/covering/identchk.rs}, $4\le m\le90$, which also fails the second known-false control $a+b\le 2t-1$ in $1601$ instances, so that the verdict of $0$ failures for the identity is not vacuous), and the argument does not use it.

Theorem~\ref{thm:pairident} subsumes Theorem~\ref{thm:pairupper}: at $q = a+b$ with $ap \equiv 1$
the residue of $b$ is $-1$, so $u = v = 1$ and the identity reads $a + b = q$. The identity is the
classical relation between two consecutive best approximations, and it holds.

\begin{theorem}[Verified]\label{thm:identity}
With the notation above, $av + bu = q$, and $\gcd(a,b) = 1$.
\end{theorem}

\begin{proof}
Let $\Lambda = \{(x,s) \in \mathbb{Z}^2 : s \equiv xp \pmod q\}$, a subgroup of index $q$
containing $(a,u)$ and $(b,-v)$, and let $\Lambda' = \langle (a,u),(b,-v)\rangle$. The determinant
of the second is $av+bu$, so $q \mid av+bu$; write $av + bu = cq$, so that $c \ge 1$ is the index of
$\Lambda'$ in $\Lambda$.

Suppose some $(z_0,s_0) \in \Lambda$ lies outside $\Lambda'$. Subtracting an element of $\Lambda'$
puts it in the fundamental parallelogram: $z = \lambda a + \mu b$ and $s = \lambda u - \mu v$ with
$\lambda, \mu \in [0,1)$ not both $0$. Then $0 < z < a+b$ and $-v < s < u$, so
$|s| < \max(u,v) < q/t$. If $z > 2m$, replace the point by $(a,u)+(b,-v)-(z,s)$, which is again
outside $\Lambda'$, has first coordinate $a+b-z$, and since $a+b \le 4m-1$ and $z > 2m$ that is
below $2m$. So we may take $1 \le z \le 2m$. Then $z$ is bad, hence $z \in \{a,b\}$. If $z = a$
then $q$ divides both $ap-u$ and $ap-s$, hence $u-s$, which lies strictly between $0$ and
$u+v < q$: impossible. If $z = b$ the same applies to $v+s$. So no such $(z_0,s_0)$ exists, that is
$\Lambda' = \Lambda$.

Two points of $\Lambda$ now give the two conclusions. First, $(1,p) \in \Lambda = \Lambda'$, so
$1$ is an integer combination of $a$ and $b$ and $\gcd(a,b) = 1$. Second, $(0,q) \in \Lambda$, and
$\Lambda' \cap (\{0\}\times\mathbb{Z}) = cq\mathbb{Z}$ because $\gcd(a,b)=1$; so $cq \mid q$,
that is $c = 1$.
\end{proof}

The construction is explicit rather than abstract: for $(0,q)$ the reduced point is
$z = (\beta a + \alpha b)/c$ and $s = (\beta u - \alpha v)/c$ with $\beta = b \bmod c$ and
$\alpha = -a \bmod c$, and both coordinates are integers by congruence alone, since
$\beta a + \alpha b \equiv ba-ab = 0$ and $\beta u - \alpha v \equiv bu+av = cq \equiv 0$ mod
$c$. For a general starting point the same formulas apply with $c$ replaced by $cq$ and
$\beta, \alpha$ the residues of $z_0v + bs_0$ and $uz_0 - as_0$. That is what is formalized.

\begin{corollary}[Verified]\label{cor:pairupper-uncond}
Every non-doubling bad pair satisfies $a+b \le 2m+2$, with no further hypothesis.
\end{corollary}

Together with Corollary~\ref{cor:pairstruct}, which gives $a+b > 2m$ in the non-doubling case, this
pins the complementary shape exactly: $a+b \in \{2m+1, 2m+2\}$. The coprimality that the shape
alone does not supply is the second half of Theorem~\ref{thm:identity}, since
$\gcd(a,a+b) = \gcd(a,b)$. The normalization asked about in earlier versions of this problem, that
every non-doubling bad pair be realized at the modulus $a+b$, is no longer needed for either.

The identity gives the bound, and the coprimality comes with it. A bad pair may have
$a+b = 2m+1$ with $a$ sharing a factor with the modulus; over $4 \le m \le 40$ this happens $25$
times, and every one of those pairs is a doubling pair, carried by the first disjunct. (The count $25$ was obtained in a session; the generating program is not archived and the count was not regenerated; HEURISTIC.) For
non-doubling pairs it does not happen, and that is a consequence rather than an observation: a time
realizing the pair at the modulus $a+b$ has $ap \equiv 1$, which exhibits $a$ as a unit. So the
normalization yields the families of Section~\ref{sec:vocab} with their hypotheses, not merely the
shapes, and that implication is formalized too.

Nothing else in this paper depends on the triple classification being exhaustive. The decision of
Theorem~\ref{thm:main-63} uses the enumeration, not the families. Theorem~\ref{thm:fam63} uses the
families, but only in the direction where omitting bad sets is harmless: a smaller family of
constraints can only lower the hitting number, so a verdict $\tau > m$ obtained from a subfamily is
a verdict for the whole family. The unclassified triples are constraints that were simply not used.

The individual families are verified. We state them with their exact hypotheses; in each case the
conclusion is an equality of sets, since a part of a cover must contain a whole bad set.

\begin{proposition}[Verified: singletons]
$\{1\}$ is the bad set of the time $1/(2m+2)$, and $\{m+1\}$ is the bad set of the time $1/(m+1)$.
\end{proposition}

\begin{proposition}[Verified: MUL]\label{prop:mul}
Let $2m/3 < d \le m$. Then the bad set of the time $1/d$ inside $[1,2m]$ is exactly $\{d, 2d\}$.
\end{proposition}

\begin{proof}
The band at modulus $d$ has half-width $d/(m+1) < 1$, so $x$ is bad exactly when $d \mid x$; the
hypothesis $3d > 2m$ says the third multiple has left the window.
\end{proof}

The number of admissible $d$ is $m - \lfloor 2m/3 \rfloor = \lceil m/3 \rceil$, and this matches the
enumerated count of MUL pairs at every $m$ in the computed range.

\begin{proposition}[Verified: CMP]\label{prop:cmp}
Let $q \in \{2m+1,\ 2m+2\}$ and let $a p \equiv 1 \pmod q$ with $2 \le a \le 2m$. Then the bad set
of $p/q$ inside $[1,2m]$ is exactly $\{a,\ q-a\}$. Moreover no other modulus supports a
complementary pair: if $\{a, q-a\}$ is a bad set with residues $\pm 1$, then $q > 2m$, since
otherwise $q$ itself is bad at residue $0$, and $q \le 2m+2$, since otherwise residue $\pm 2$ is bad
and the set is larger.
\end{proposition}

The case $q = 2m+1$ is the coprimality edge of Bohman and Peng \cite[\S3]{BP}, who take the time
$s^{-1}/(s+t)$ and observe that no velocity is $0$ or $\pm 1$ modulo $s+t$; we cite it and do not
claim it. The case $q = 2m+2$ appears not to be recorded. At $a = 1$ the partner $q-1 = 2m+1$ leaves
the window and the set degenerates to the singleton $\{1\}$, consistently with the singleton family.

\begin{proposition}[Verified: DBL]\label{prop:dbl}
Let $2m+2 < q \le 3m+3$, let $ap \equiv -1 \pmod q$, and suppose $2a \le 2m$ and $2m + 2a < q$. Then
the bad set of $p/q$ inside $[1,2m]$ is exactly $\{a, 2a\}$.
\end{proposition}

The hypothesis $2m + 2a < q$ pushes the complementary partners $q-a$ and $q-2a$ out of the window;
without it they re-enter and the bad set has four elements.

The classification of pairs has a structural consequence for the graph they form. Write $G_m$ for
the graph on $[2,2m]$, on $2m-1$ vertices, in which $a$ and $b$ are adjacent when $\{a,b\}$ is a bad
set. The vertex $1$ is omitted because it is a bad singleton. Section~\ref{sec:obstructions} works
instead with $G$, the induced subgraph of $G_m$ on $[2,2m]\setminus\{m+1\}$, on $2m-2$ vertices,
which drops the second bad singleton as well. The two are distinct and both are used: triangle
freeness is stated for the larger $G_m$, and the matching and independence results for $G$, where
the candidate complements live.

\begin{theorem}[Verified]\label{thm:trianglefree}
$G_m$ contains no triangle, for every $m$.
\end{theorem}

\begin{proof}
Every edge is a doubling pair $b = 2a$ or a complementary pair with $a + b \in \{2m+1, 2m+2\}$, by
Corollary~\ref{cor:pairstruct} and Theorem~\ref{thm:pairupper}. Let $2 \le x < y < z$ be a
triangle. If all three edges are complementary, the three sums lie in $\{2m+1,2m+2\}$, which forces
$x,y,z \in [m-1,m+2]$; the only triple there with $x+y \ge 2m+1$ is $\{m,m+1,m+2\}$, and it fails
because $y+z = 2m+3$. If exactly one edge doubles, subtracting the two complementary sums gives
$x \le 1$. If two edges double, the triangle is $\{x,2x,4x\}$ with $5x \in \{2m+1,2m+2\}$, and the
gcd condition then gives $x = 1$. Three doubling edges are inconsistent with $x<y<z$. Every branch
contradicts $x \ge 2$.
\end{proof}

Two consequences are used in Section~\ref{sec:obstructions}. A part of three elements and capacity
one requires all three of its pairs to be bad, so no such part exists. And no five-cycle of $G_m$ has
all five edges complementary, since two complementary edges at a vertex force its neighbours to
differ by at most one, which places all five vertices in a window of width two.

For triples we record the three families and omit the proofs, which are of the same character.

\begin{proposition}[Verified: SUM, SUM$'$, SUM$''$, DBL$_3$]\label{prop:triples}
Let $a < b$.
\begin{enumerate}
\item If $m+1 < a+b \le 2m$, $2a+b > 2m$ and $a p \equiv 1 \pmod{a+b}$, the bad set of $p/(a+b)$ is
  exactly $\{a, b, a+b\}$.
\item If $b \ge m+1$, $a + b \le 2m$, $a + 2b \le 3m+3$ and $b p \equiv 1 \pmod{a+2b}$, the bad set
  of $p/(a+2b)$ is exactly $\{a, b, a+b\}$.
\item If $2a+b \le 2m$, $b \le m$, $a+2b > 2m$, $m+1 < a+b$, $3m+3 < 2a+3b \le 4m+4$ and
  $(a+b)p \equiv 1 \pmod{2a+3b}$, the bad set of $p/(2a+3b)$ is exactly $\{a, b, a+b\}$. The
  residues of $a$, $b$, $a+b$ at this modulus are $3$, $-2$, $1$.
\item If $2a \le 2m$, $t \le 2m$, $a + t > 2m$, $2m+2 < 2a+t \le 3m+3$ and $ap \equiv 1
  \pmod{2a+t}$, the bad set of $p/(2a+t)$ is exactly $\{a, 2a, t\}$.
\end{enumerate}
\end{proposition}

The third item is family SUM$''$. It is the branch the enumeration of
Section~\ref{sec:computation} has always used and this proposition did not state; the lower bound
$3m+3 < 2a+3b$ is necessary and was absent from the enumeration's predicate. The smallest member
carries residue $3$, so it is bad only when $3(m+1) < 2a+3b$, while the remaining conditions give
only $2a+3b > 3m+1$. Without it the branch admits sets whose smallest member is not bad, of which
there are $38$ below $m = 60$, the first being $\{3,4,7\}$ at $m = 5$, where $2a+3b = 18$ equals
$3(m+1)$ and the bad set is $\{4,7\}$. With it the branch is realizable throughout the tested
range, and the enumeration's membership check is unaffected.

\begin{remark}[Formalization]
The side condition $a+t>2m$ of item (3) need not be stipulated. Write $\theta_x$ for the signed
distance from $x\alpha$ to the nearest integer, so that $x$ is bad exactly when
$|\theta_x|<\varepsilon$, and note that on the window the residues add. For an exact bad set
$\{a,2a,t\}$ the sum $a+t$ exceeds $t$, so if it were bad and at most $2m$ it would be a member
larger than the largest, which is impossible; hence either $a+t>2m$, or $a+t$ is not bad and the
difference $t-a$ is, and $t-a$ lies strictly below $t$, so it is $a$ or $2a$ and $t=3a$. Thus
$a+t\le 2m$ forces the chain $\{a,2a,3a\}$, and item (3), where $t\neq 3a$, must satisfy $a+t>2m$.

The same argument gives the shape dichotomy: an exact bad set $\{a,b,c\}$ with $a+b\le 2m$ satisfies
$c=a+b$ or $b=2a$. For $a+b>2m$ the sum branch is vacuous, since badness outside $[1,2m]$ identifies
no member, and the reflection of Theorem~\ref{thm:identity} is required: a reduced lattice point
$z=\lambda a+\mu c$ above the window reflects to $a+c-z\in[1,2m]$ with residue
$(1-\lambda)\theta_a+(1-\mu)\theta_c$, and since $z>2m$ forces $\lambda+\mu>1$, that residue is
bounded by $(2-\lambda-\mu)\max(|\theta_a|,|\theta_c|)$ and is therefore strictly smaller. The
reflected point is bad and in range, hence a member.

The determinant argument behind Lemma 3.1 is formalized in the same way. With
$\theta_a=a\alpha-j_a$ and $\theta_b=b\alpha-j_b$ the quantity $a\theta_b-b\theta_a$ equals
$b j_a-a j_b$, an integer, the $\alpha$ terms cancelling; and it is bounded by $(a+b)\varepsilon$,
which stays below $2$ on the window. An integer of absolute value below $2$ is $-1$, $0$ or $1$, and
the vanishing case makes the two lifts parallel, so that $\theta_x=x\delta$ with $\delta=\alpha-k$
and every $x\le a$ is bad, forcing the least bad speed to be $1$. Hence the determinant is $\pm1$,
which is itself the B\'ezout relation: $b j_a-a j_b=1$ reads $a(-j_b)+b j_a=1$, so $\gcd(a,b)=1$ and
the coefficients are exhibited at once.

These statements are formalized in Lean, without \texttt{sorry} and on the standard axioms, in
\texttt{ShapeTheorem.lean}, \texttt{ReflectionResidue.lean} and \texttt{LatticeReduction.lean}. The membership of every bad set of
size at most three in the five families is checked for $2\le m\le 150$, over $q\le 2m+2$, by kernel evaluation and for
$2\le m\le 300$, comprising $5{,}882{,}334$ sets, by an independent exhaustive sweep. The shape
dichotomy is proved for $a+b\le 2m$. The complementary case is equivalent to a sharper statement: an
exact bad set whose two smallest members form a DBL pair with $3a>2m$ has no third member. Over
$2\le m\le 150$ there are $194{,}811$ such pairs and none extends (program \texttt{code/covering/bt\_dbl\_extend.rs}), so the statement is non-vacuous,
but it is not proved. It cannot be proved from the residue inequalities alone: with $a>2m/3$ every
speed is $r$, $a+r$ or $2a+r$ with $0\le r<a$, a third member exists exactly when some $r$ satisfies
$|2\theta_a+\theta_r|<\varepsilon$ while $r$ and $a+r$ are not bad, and that system is satisfiable,
for instance at $\theta_a=\varepsilon/4$, $\theta_r=-5\varepsilon/4$. The obstruction must come from
which residues occur among the available $r$, which is a three-distance question.
\end{remark}

Families (1) and (2) realize the same triple $\{a,b,a+b\}$ at two different moduli; (1) applies when
the second periods stay outside the window and (2) when they do not. The distinction is not
cosmetic: at $m = 7$ the triple $\{2,9,11\}$ fails the hypothesis of (1), since $2a+b = 13 \le 14$,
and the bad set at modulus $11$ is in fact $\{2,9,11,13\}$; it is realized by (2) at modulus $20$.

\begin{remark}
An empirical observation, computed and not proved, guided the choice of families. Extracting an
irredundant subfamily of bad sets that still forces $\tau(2m) > m$, by repeatedly discarding a bad
set and re-deciding, produces for every $m \le 33$ a core containing no bad set of size four or
more. When $2m+1$ is prime the core has exactly $m+1$ members, namely the two singletons and $m-1$
pairs forming a perfect matching, which recovers the structure behind the uniform proof for prime
$2m+1$. When $2m+1$ is composite the core is larger and always contains triples. This is why the
classification stops at size three.
\end{remark}

\section{Decision procedure and computations}\label{sec:computation}

Theorem~\ref{thm:equiv} turns the window statement into the decision problem: does a hitting set of
size $m$ exist? We decide it by branch and bound, seeding the incumbent at $m+1$, branching on the
smallest not-yet-hit bad set, and pruning with the bound that pairwise disjoint not-yet-hit bad sets
each require a distinct element. Bad sets are enumerated exhaustively as in Section~\ref{sec:vocab}
and reduced to minimal ones.

A verdict of the form ``no hitting set of size $m$ exists'' is robust to under-enumeration: $\tau$
is monotone in the family of sets to be hit, so establishing $\tau > m$ for a subfamily establishes
it for the whole family. Only a verdict in the other direction would require completeness of the
enumeration, and none occurred.

\begin{proof}[Proof of Theorem~\ref{thm:main-63}]
The procedure was run for $4 \le m \le 63$, the ceiling being the width of the machine word used for
the window bitmask. In every case the search exhausted without finding a hitting set of size $m$,
whence $\tau(2m) = m+1$ by Theorem~\ref{thm:ceiling} and the window statement follows by
Theorem~\ref{thm:equiv}. Total running time was under thirteen minutes, dominated by the enumeration
of bad sets rather than by the search.
\end{proof}

\subsection*{The same range from the verified families alone}
The monotonicity just used has a consequence worth separating out. Since a verdict $\tau > m$ for a
subfamily already gives the window statement, the search does not need the exhaustive enumeration at
all: it may be run against the nine families of Section~\ref{sec:vocab} and nothing else. Those
families are verified, so on this route the enumeration, which is the one ingredient of
Theorem~\ref{thm:main-63} that is neither formalized nor independently auditable, is removed
entirely.

\begin{theorem}[Computed, from verified inputs]\label{thm:fam63}
For every $m$ with $4 \le m \le 63$, the hitting-set search run against the two singleton, four pair
and three triple families of Section~\ref{sec:vocab}, and against no other set, exhausts without
finding a hitting set of size $m$. Consequently the window statement holds for every such $m$.
\end{theorem}

The soundness of this route depends on an inclusion in one direction only, and it is the opposite of
the one in the previous paragraph. Under-generation is harmless. Over-generation is not: a set that
is not a bad set is a constraint with no content behind it, and since adding constraints can only
raise $\tau$, a single such set could manufacture a false verdict $\tau > m$. What is required is
therefore exactly that every generated set is a bad set, and this is certified constructively rather
than assumed. Each family member is emitted together with the time $p/q$ named in its proposition,
and $B_{2m}(p,q)$ is recomputed from the definition and compared with the generated set. Over
$4 \le m \le 250$ this was carried out for all $2{,}772{,}478$ generated members; in every case the
two sets were equal, not merely nested. The certificate needs no enumeration and is therefore
available at every $m$, including those far beyond the range the search itself reaches.

Four controls were run. First, the verdicts for $m \le 24$ agree with the independently verified
range. Second, replacing the threshold $1/(m+1)$ by $1/m$ makes the statement false, since
$\ML(1,\dots,m) = 1/(m+1)$; the procedure then finds a hitting set, and the witness it returns is
$\{1,\dots,m\}$, the classical tight instance. Third, an independent implementation computing
$\tau(2m)$ exactly, with a different bound and code path, returns $\tau = m+1$ at $m = 25$ and
$m = 26$, in agreement. Fourth, a search in the complementary formulation, looking for an $m$-set
containing no bad set rather than for a hitting set, returns the same verdict at every $m \le 45$;
it shares no branching logic with the first implementation.

\begin{remark}\label{rem:sweep}
The enumeration of bad sets can be made much cheaper than the scan over all reduced $p/q$ with
$q \le 2N(m+1)$, which costs $O(m^4)$ times the window width. Membership of $x$ changes only where
$\Norm{x\alpha} = 1/(m+1)$, that is at $\alpha = (j(m+1) \pm 1)/(x(m+1))$, so there are only $O(m^2)$
breakpoints in $[0,1)$; evaluating the bad set at every breakpoint and at one interior point of every
interval between consecutive breakpoints therefore produces the complete family in $O(m^3)$. The
breakpoints must be evaluated and not only the intervals: at a breakpoint some $x$ has
$\Norm{x\alpha} = 1/(m+1)$ exactly and is good there while bad on both sides, so the bad set at a
breakpoint is a proper subset of both its neighbours, and a smaller bad set is a stronger
constraint. Carried out in exact integer arithmetic on rationals, with interior points taken as
mediants, this reproduces the family of minimal bad sets produced by the scan exactly, set for set,
for every $m \le 26$. With the enumeration no longer the bottleneck, the branch and bound is, and its
cost tracks the deficiency rather than $m$, which is what Proposition~\ref{prop:ke} would lead one
to expect: searching against the nine families, $m = 56$ takes under a second and $m = 61$, whose
deficiency is $18$, takes half an hour.
\end{remark}

\begin{remark}\label{rem:lift}
Theorem~\ref{thm:main-63} is computed, not verified. The decision procedure establishes the
nonexistence of a small hitting set, which is not by itself a certificate that the Lean statement
can consume; the verified entry points require a cover with capacities. By Lemma~\ref{thm:surplus}
and Proposition~\ref{prop:surplus-one}, producing such a cover at a given $m$ amounts to
partitioning $[1,2m] \setminus \{1,m+1\}$ into parts of surplus zero, and we do not know how to do
this uniformly; see Section~\ref{sec:obstructions}.
\end{remark}

\subsection*{A case above the window: $m = 13$ at height three}

The decision procedure is not tied to the window $[1,2m]$. Raising the height to $3m$ at $m = 13$
gives a case the literature does not cover, since the conjecture is known for at most $12$
velocities and that is independent of height. We record the convention once more because a single
velocity separates new from known: $m$ counts velocities and the number of runners is $m+1$, so
$m = 13$ is $14$ runners.

\begin{theorem}[Computed]\label{thm:k3m13}
Every $13$-element set of velocities inside $[1,39]$ satisfies \eqref{eq:lrc}.
\end{theorem}

The criterion is the hitting formulation of Section~\ref{sec:hitting} applied verbatim at height
three. A velocity set is lonely exactly when some time's bad set misses it entirely, so a
counterexample is a set meeting every bad set. Over $[1,39]$ at threshold $1/14$ the family of
distinct bad sets has $717$ members, obtained from the breakpoints together with the reduced
mediants between consecutive ones. All $\binom{39}{13} = 8\,122\,425\,444$ velocity sets were
tested against it and none meets every member.

The direction of any error in the enumeration is the one that matters. Each member of the family is
computed from an explicit rational by the definition, so the family can only omit bad sets, never
contain a set that is not bad. An omission would report a lonely set as a counterexample, not the
reverse, so a count of zero is safe while a positive count would have required checking each witness
by hand.

Two independent checks support the computation. First, the same procedure run at $m = 10, 11, 12$
and the same height $3m$ returns no counterexample in each case, agreeing with the literature, which
settles those by other methods. Second, the procedure detects counterexamples when they exist: at
threshold $1/13$ rather than $1/14$ the set $\{1,\dots,13\}$, whose maximum loneliness is exactly
$1/14$, is reported, and at threshold $1/10$ in $[1,20]$ four sets are reported. An earlier version
of the program discarded empty bad sets, which are the times at which nothing in the window is bad,
and consequently reported $\{1,\dots,13\}$ as a counterexample at threshold $1/14$; the second check
is what caught it.

Theorem~\ref{thm:k3m13} is computed and not verified, and the obstruction is sharper than the one in
Remark~\ref{rem:lift}. At height $b+1$ the budget of Lemma~\ref{thm:surplus} becomes $b(k+1)$, so
at $m = 13$ and height three a certificate needs parts covering $[1,39]$ with capacities summing to
at most $25$, an average ratio of $0.641$. At this height there is no bad singleton and there are
only two bad pairs, so the parts available at or below that ratio are scarce: one of size four, $54$
of size five, none of size six with capacity three, and eight of size seven, the last all containing
$\{1,2,14,28\}$. Every element of $[1,39]$ lies in some such part, so nothing is excluded outright,
yet an exhaustive search over parts of size two through seven finds no cover within budget. The
scarce elements, each lying in exactly three candidate parts, are what the search fails on.

\section{Closed approaches}\label{sec:obstructions}

We record five approaches to the uniform statement that are now closed, with reasons. Each was
pursued to a decision, and the negative results are the durable part of this work.

\subsection*{Bad pairs alone, integrally or fractionally}
A cover by the two singletons and $m-1$ disjoint bad pairs certifies the window, and such a matching
exists precisely when $2m+1$ is prime or equal to $9$, which was checked by blossom matching for
every $m \le 24$. The fractional relaxation does no better: by half-integrality of the fractional
matching polytope a vertex set decomposes into edges and odd cycles exactly when the graph has a
fractional perfect matching, and computation over $5 \le m \le 24$ shows the bad-pair graph has one
if and only if $2m+1$ is prime.

Let $G$ be the bad-pair graph on $[1,2m] \setminus \{1,m+1\}$, on $n = 2m-2$ vertices, that is the
induced subgraph of $G_m$ obtained by deleting $m+1$, with
Gallai--Edmonds decomposition $(D,A,C)$ and deficiency $\mathrm{def} = n - 2\nu(G)$.

\begin{proposition}[Proved; Classical]\label{prop:alpha}
$\alpha(G) \le m - 1 + \tfrac12 \mathrm{def}$, with equality if and only if $G$ is
K\"onig--Egerv\'ary, that is $\alpha(G) + \nu(G) = n$.
\end{proposition}

We claim nothing here: the statement is Gallai's identity combined with the definition of the
K\"onig--Egerv\'ary property, and we record it only to fix the direction of the inequality, which
is the direction that makes it useless on its own.

\begin{proof}
Gallai's identity $\alpha(G) + \rho(G) = n$, with $\rho$ the minimum vertex cover, together with
$\rho(G) \ge \nu(G)$, gives $\alpha(G) \le n - \nu(G)$, and equality holds exactly when
$\rho(G) = \nu(G)$, which is the K\"onig--Egerv\'ary property. Substituting
$\nu(G) = (n - \mathrm{def})/2$ and $n = 2m-2$ turns $n - \nu(G)$ into $m - 1 + \tfrac12\mathrm{def}$.
\end{proof}

The inequality runs in the direction that is useless on its own: it caps $\alpha(G)$ rather than
forcing it up. What the computation supplies is the equality case.

\begin{proposition}[Computed]\label{prop:ke}
For every composite $2m+1$ with $m \le 63$ the bad-pair graph is K\"onig--Egerv\'ary, so
$\alpha(G) = m - 1 + \tfrac12\mathrm{def}$ exactly. The independence number was computed by exact
maximum independent set, by a code path that does not use the deficiency, and the deficiency by
blossom matching; the identity $\alpha + \nu = n$ was then read off. Consequently
$\alpha(G) \ge m$ at each of the thirty-two of those values with $\mathrm{def} \ge 2$, the one
exception being $m = 4$, where $2m+1 = 9$ and $\mathrm{def} = 0$.
\end{proposition}

Suppose $\alpha(G) \ge m$ and let $S$ be an independent set of $G$ with $|S| = m$. Its complement
$H = [1,2m] \setminus S$ also has $m$ elements, and it meets every bad set of size at most two:
it contains $1$ and $m+1$, which lie outside the ground set of $G$ and so outside $S$, and it meets
every bad pair, since a bad pair inside $S$ would be an edge of $G$ inside an independent set. Hence
the hitting number of the subfamily consisting of the bad singletons and the bad pairs is at most
$m$, and no certificate assembled from that subfamily can reach $m+1$. So at every composite $2m+1$
with $m \le 63$ other than $m = 4$, no certificate built from bad pairs alone can work, integrally
or fractionally, and the bad triples are doing work that nothing smaller can do.

This is a computed statement, and it cannot be promoted, because the K\"onig--Egerv\'ary property
is \emph{false} in general.

\begin{proposition}[Computed]\label{prop:ke-false}
The bad-pair graph is not K\"onig--Egerv\'ary for every $m$. Over $4 \le m \le 900$ it fails at
exactly four values, $m = 150, 318, 486, 570$. The least is $m = 150$, where $2m+1 = 301 = 7\cdot43$:
there $|V| = 298$, $|E| = 323$ and $\nu = 139$, and no vertex cover of size $\nu$ exists.
\end{proposition}

All four have $7 \mid 2m+1$. We state that as an observation and not as a pattern: it is four points,
the divisibility is necessary over the tested range but plainly not sufficient, since $128$ of the
$897$ values tested satisfy it and only four fail, and the tightest residue class containing all
four still contains six values that do not fail. Whether the failures are finite, sparse, or merely
rare is open, and nothing in this paper depends on the answer.

Whether a graph is K\"onig--Egerv\'ary is decidable in polynomial time, without computing the
independence number, which is what makes a search out to $m = 150$ possible. Fix a maximum matching
$M$. A vertex cover of size $|M|$ must contain exactly one endpoint of each edge of $M$ and no
$M$-exposed vertex, since the $|M|$ disjoint edges already exhaust the budget; one boolean per edge
of $M$ then turns the covering requirement into a $2$-SAT instance, satisfiable exactly when the
graph is K\"onig--Egerv\'ary. In the satisfiable case the assignment is an explicit minimum vertex
cover, so a positive answer certifies itself. K\"onig--Egerv\'ary graphs carry a literature of
their own \cite{Dem,Ste,Lar}; Proposition~\ref{prop:ke-false} says the bad-pair graph does not
belong to it, so that literature is not a route to a proof here.

The equality $\alpha(G) = m - 1 + \tfrac12\mathrm{def}$ is therefore a feature of small $m$ rather
than a theorem, and the negative result above rests on the measured values of $\alpha(G)$ and on
nothing else. One half of it is unconditional and worth isolating, since it needs neither
computation nor the classification.

\begin{proposition}[Verified]\label{prop:alpha-m-1}
$\alpha(G) \ge m-1$ for every $m$.
\end{proposition}

\begin{proof}
By Theorem~\ref{thm:ceiling} every bad set meets $[1,m+1]$. Hence no bad set is contained in
$[m+2,2m]$, so that interval contains no bad set at all; in particular it spans no edge of $G$, and
it is an independent set of $G$ of size $m-1$.
\end{proof}

The three steps are formalized: that the bad set always escapes the top interval, that no bad set is
contained in it, and its specialization to the two-element sets which are the edges of $G$, together
with the cardinality $m-1$.

The remaining unit, from $m-1$ to $m$, reduces to a statement about a single element. Write
$T = [m+2,2m]$ and let $x \notin T$. Since no bad set is contained in $T$, every bad set contained in
$T \cup \{x\}$ must contain $x$; this too is verified. So $T \cup \{x\}$, which has $m$ elements, is
independent in $G$ as soon as $\{x\}$ is not a bad set and no $\{x,y\}$ with $y \in T$ is a bad set.
Nothing about how bad sets distribute over $T$ enters.

\begin{proposition}[Verified, given the classification]\label{prop:topgap}
Let $2m+1$ be composite and let $p$ be its least prime factor. Assume the classification of
Theorem~\ref{thm:oth} in the form that a doubling bad pair $\{a,2a\}$ satisfies the hypothesis of
MUL or of DBL. Then for every $m \ge 13$ the element $x = 2p$ forms no bad pair with any element of
$T$, so $\alpha(G) \ge m$. The exceptions below that range are $m = 7$ and $m = 12$, where some
element other than $2p$ is needed, and $m = 4$, where $2m+1 = 9$ and no element works at all.
\end{proposition}

\begin{proof}
Since $p^2 \le 2m+1$ we have $2p < m+2$, so $2p$ lies below $T$ and the pair is $\{2p, y\}$ with
$y \in T$. If it is complementary at $2m+1$ then $a$ is a unit there, but $p$ divides both $2p$ and
$2m+1$; if complementary at $2m+2$ then $2$ divides both. If it is a doubling pair then $y = 4p$.
Under MUL, $3\cdot 2p > 2m$ gives $3p > m$, so $9p^2 > m^2$; with $p^2 \le 2m+1$ this forces
$m^2 < 18m+9$, hence $m \le 18$. Combined with $m \ge 13$ and $3p > m$ this gives $p \ge 5$, while
$p^2 \le 2m+1 \le 37$ gives $p \le 6$, so $p = 5$; then $3p > m$ forces $m \le 14$ and $5 \mid 2m+1$
fails, since $2m+1$ is $27$ or $29$. Directly: for each $m$ with $13 \le m \le 18$ and $2m+1$
composite, namely $m \in \{13,16,17\}$ with least prime factor $3,3,5$, the condition $3p > m$
already fails. Under DBL
the modulus condition $2m + 2\cdot 2p < 3m+3$ gives $4p < m+3$, while $y = 4p \in T$ needs
$4p \ge m+2$; so $4p = m+2$, whence $2m+1 = 8p-3$ and $p \mid 3$, giving $p = 3$ and $m = 10$.
Both are below $13$.
\end{proof}

Only $m = 17$ and $m = 24$ are threatened at all in $13 \le m \le 28$, and at both the MUL and DBL
hypotheses fail. The weaker form of the classification, shape together with coprimality, suffices
for $m \ge 29$, where the doubling case is excluded by size alone.

The three failures are of two different kinds, and the distinction matters. At $m = 4$ the
deficiency is $0$ and $\alpha(G) = m-1$ exactly, so no element can work and nothing is missing. For
prime $2m+1$ no element works either, again correctly, and for the same reason. At $m = 7$ and
$m = 12$, by contrast, $\alpha(G) \ge m$ does hold while no set of the shape $T \cup \{x\}$ attains
it, so there the construction is genuinely insufficient rather than merely unlucky.

\begin{proposition}[Verified]\label{prop:topgap-fails}
For $m = 7$ and for $m = 12$, every $x \in [2,m]$ admits a time whose bad set inside the window is a
two-element subset of $T \cup \{x\}$. Hence no set of that shape is independent in $G$ at those two
values.
\end{proposition}

Each of the seventeen blocking pairs is certified by an explicit time, and membership in a bad set
is decidable, so the proposition is settled case by case inside the proof assistant rather than read
off a table. We formalize it because the shape that fails is as much a part of the record as the
shape that works.

Why $2p$ should work is visible from the families of Section~\ref{sec:vocab}: $2p$ is a non-unit
modulo $2m+1$, so it carries no CMP$(2m+1)$ pair; it is even while every unit modulo $2m+2$ is odd,
so it carries no CMP$(2m+2)$ pair; and its doubling partners $p$ and $4p$ are both below $m+2$ once
$p \le \sqrt{2m+1}$, so the MUL and DBL pairs through it do not reach $T$. That argument is not a
proof, because it presumes the classification exhaustive, which is Problem~2 of
Section~\ref{sec:open}.

\subsection*{Enumerating the independent sets}
Proposition~\ref{prop:ke} makes the residual task explicit: every independent $m$-set of $G$ must
contain a bad set of size at least three. Writing $I(m)$ for the number of independent $m$-sets, the
computed values (program \texttt{code/covering/indcount.rs}, $4\le m\le24$; it returns $I(m)=0$ at every prime $2m+1$ in the range) show $I(m)$ is governed by the deficiency rather than by $m$:
$I = 1, 9, 3, 41, 1203$ at $\mathrm{def} = 2$ (for $m = 7,10,12,13,24$), $I = 310, 728$ at
$\mathrm{def} = 4$, $I = 4093$ at $\mathrm{def} = 6$, and $I = 40493$ at $\mathrm{def} = 8$. Note
$m = 22$, with $\mathrm{def} = 8$, has $I = 40493$ while the larger $m = 24$, with
$\mathrm{def} = 2$, has $I = 1203$.

Over the computed values $I$ grows by roughly an order of magnitude per two units of deficiency, so
that $I \approx 10^{\mathrm{def}/2}$ describes the data. If the deficiency is $\Theta(m)$ for
infinitely many $m$, which we expect but do not prove, and if that empirical growth persists, then
$I(m)$ is exponential in $m$ along those $m$, and any argument treating the independent $m$-sets
individually, by enumeration or by case analysis, is not uniform in $m$. We state this as the
heuristic it is; it is a guide to the shape of an admissible proof rather than a proved obstruction,
and each of the approaches below violates it.

\subsection*{Disjoint bad sets at greater height} By Proposition~\ref{prop:packing} a certificate of
pairwise disjoint bad sets inside $[1,Bm]$ needs $B$ bad sets of the minimum size $B-1$, while by
Proposition~\ref{prop:surplus-two} only two exist. Hence this shape of certificate is confined to
$B = 2$.

\subsection*{Prime square moduli} The Gallai--Edmonds decomposition of the bad-pair graph satisfies
$|D| - |A| = \operatorname{def}$ at all thirty-three composite $2m+1$ with $m \le 63$, so over that
range $D$ is an independent set and the obstruction is a Hall violator rather than an odd component.
We state the range because it is the whole evidence: the K\"onig--Egerv\'ary regularity of
Proposition~\ref{prop:ke} was measured over exactly the same values and fails at $m = 150$. One hopes for an arithmetic
description of $D$. Two natural candidates fail: $D$ is not the set of non-units modulo $2m+1$, and
the deficiency is not the number of odd doubling chains among the non-units, agreeing in only seven
of fourteen computed cases and failing in both directions. For $q = 2m+1 = p^2$ the deficiency takes
the values $0, 2, 2, 4$ at $p = 3,5,7,11$, so the family is not uniform either.

\subsection*{The six-three template} Every non-overlapping certificate we verified consists of the
two singletons together with $(m-1)/3$ parts of six elements and capacity three, which requires
$3 \mid m-1$. This holds at $m = 10, 13, 16$. It fails at $m = 19$: enumerating all six-element
subsets of $[1,38] \setminus \{1,20\}$ and retaining the $4440$ with capacity exactly three, an
exhaustive exact-cover search terminates without a partition after $969{,}154{,}513$ nodes (program \texttt{code/covering/tmpl63.rs}; its default node cap is $2\cdot10^{8}$, so the cap must be raised; the same program finds the template at $m=10,13,16$). So the
template is a coincidence of small $m$, not a family.

\section{Formalization}\label{sec:lean}

The development is in Lean 4 with Mathlib. Every statement labelled verified builds cleanly with no
\texttt{sorry} and depends only on \texttt{propext}, \texttt{Classical.choice} and
\texttt{Quot.sound}; the few statements that use no choice are noted in the source. We are not aware
of a previous formalization of any part of the lonely runner conjecture in any proof assistant.

\begin{center}
\begin{tabular}{@{}ll@{}}
\toprule
Statement & Lean declaration \\
\midrule
Theorem~\ref{thm:mincount} & \texttt{minBadCount} \\
Theorem~\ref{thm:ceiling} & \texttt{bad\_meets\_initial} \\
Theorem~\ref{thm:equiv}(1) & \texttt{floor\_of\_no\_hitting\_set} \\
Theorem~\ref{thm:equiv}(2) & \texttt{not\_floor\_of\_hitting\_set} \\
converse witness & \texttt{exists\_good\_time\_of\_floor} \\
Lemma~\ref{thm:surplus} & \texttt{floor\_of\_surplus} \\
capacity certificate, height $B$ & \texttt{floor\_of\_cover\_capacity\_height} \\
Proposition~\ref{prop:packing} & \texttt{packing\_needs\_minimal} \\
Proposition~\ref{prop:alpha-m-1} & \texttt{no\_badpair\_in\_top}, \texttt{card\_top\_interval} \\
\quad supporting steps & \texttt{bad\_escapes\_top}, \texttt{no\_badset\_in\_top} \\
reduction to a single element & \texttt{badset\_in\_top\_insert}, \texttt{card\_top\_insert} \\
Proposition~\ref{prop:topgap-fails} & \texttt{top\_insert\_blocked\_m7}, \texttt{top\_insert\_blocked\_m12} \\
Theorem~\ref{thm:pairdich} & \texttt{bad\_diff\_or\_sum} \\
Corollary~\ref{cor:pairstruct} & \texttt{badpair\_doubling\_or\_large} \\
Theorem~\ref{thm:pairupper} & \texttt{pairUpper\_at\_modulus} \\
Proposition~\ref{prop:topgap} & \texttt{twoP\_closes\_gap\_full}, \texttt{twoP\_closes\_gap} \\
\quad reduction chain & \texttt{pairUpperBound\_of\_realized}, \texttt{pairClassification\_of\_realized} \\
\quad with coprimality & \texttt{pairClassificationStrong\_of\_realized} \\
Proposition~\ref{prop:mul} & \texttt{badset\_mul\_pair} \\
Proposition~\ref{prop:cmp}, $q = 2m+1$ & \texttt{badset\_unit\_pair} \\
Proposition~\ref{prop:cmp}, $q = 2m+2$ & \texttt{badset\_compl\_pair} \\
Proposition~\ref{prop:dbl} & \texttt{badset\_dbl\_pair} \\
Proposition~\ref{prop:triples}(1),(2),(3) & \texttt{badset\_sum\_triple}, \texttt{badset\_sum\_triple\_high}, \texttt{badset\_dbl\_triple} \\
all windows imply the conjecture & \texttt{lonelyRunner\_of\_all\_windows} \\
dense window, $2m+1$ prime & \texttt{dense\_floor\_of\_prime} \\
dense window, $2m+1$ composite, $m \le 24$ & \texttt{dense\_floor\_m4}, \texttt{dense\_floor\_m7}, \dots, \texttt{dense\_floor\_m24} \\
Theorem~\ref{thm:identity}, the reflection step (gap G1) & \texttt{dbl\_reflection} \\
\bottomrule
\end{tabular}
\end{center}

The programs for the computed statements are in \texttt{code/covering} and \texttt{code/dense} of the
accompanying archive; its \texttt{README.md} lists, statement by statement, which counts are regenerated by a
program in the archive and which are not.

Statements that are computed or conjectural, rather than verified, are named in Lean as predicates
and are never asserted. Naming them fixes their quantifiers and the direction of their inequalities
before any effort is spent on a proof, and a statement that cannot be typed out correctly is not yet
understood; but a proof obligation discharged by \texttt{sorry} is an assertion the kernel accepts
on faith, and no statement this paper relies on is discharged that way. Every result that needs an
open statement takes it as an
explicit hypothesis, so the dependency appears in the type rather than in a comment. The open
statements are the identity behind the pair classification, the independence statement
$\alpha(G) \ge m$, and the window on $25 \le m \le 63$.

Consequently \texttt{\#print axioms} on any declaration cited in the table above returns
\texttt{propext}, \texttt{Classical.choice} and \texttt{Quot.sound}, or a subset of them, and never
\texttt{sorryAx}. That is a mechanical check, and it is what separates what this paper establishes
from what it assumes. No declaration in the Lean modules on which the table above depends carries a \texttt{sorry} body. The four
stated-first contracts that previously did, for eight runners and for the residual case of four, are
held outside it; they were imported by nothing cited here, and they remain conjectural.

\section{Where the window statement now stands}\label{sec:reduction}

The prime case is settled outright, and we state it because the reduction below rests on it.

\begin{theorem}[Verified]\label{thm:prime}
Let $2m+1$ be prime. Then every set of $m$ distinct positive velocities bounded by $2m$ satisfies
\eqref{eq:lrc}.
\end{theorem}

\begin{proof}
The complementary pairs $\{a, 2m+1-a\}$ for $1 \le a \le m$ partition $[1,2m]$, and each satisfies
$\gcd(a, 2m+1) = 1$ because $2m+1$ is prime, so each is a bad pair by
Proposition~\ref{prop:cmp}. A complement $W$ of size $m$ containing no bad set is independent in
the bad-pair graph, so it meets each of those $m$ pairs at most once, hence exactly once. It avoids
$1$ and $m+1$, both bad singletons, so from the pair $\{1,2m\}$ it takes $2m$ and from the pair
$\{m,m+1\}$ it takes $m$. But $\{m,2m\}$ is a bad pair, by Proposition~\ref{prop:mul} at $d = m$.
So no such $W$ exists, and by Theorem~\ref{thm:equiv} the window statement holds.
\end{proof}


The classification of Section~\ref{sec:vocab} reduces the window statement to a single unproved
combinatorial claim. We record the reduction because it isolates what remains.

Suppose $V \subseteq [1,2m]$ has $m$ elements and defeats \eqref{eq:lrc}. By
Theorem~\ref{thm:equiv} its complement $W = [1,2m] \setminus V$, also of size $m$, contains no bad
set. Three consequences are immediate and each is verified.

\begin{enumerate}
\item $W$ contains neither $1$ nor $m+1$, since both are bad singletons.
\item $W$ contains no bad pair, so $W$ is independent in the bad-pair graph.
\item If $2m+1$ is prime the bad-pair graph has a perfect matching on
  $[1,2m] \setminus \{1,m+1\}$, and the independence number is below $m$; this is
  Theorem~\ref{thm:prime}, and no such $W$ exists.
\end{enumerate}

So only composite $2m+1$ remains, and there the independence number is at least $m$: pairs alone
cannot finish, and the bad triples must. What is missing is exactly this.

\begin{conjecture}[Conjecture]\label{conj:link}
Let $2m+1$ be composite. Every $m$-element subset of $[1,2m]$ that avoids $1$ and $m+1$ and is
independent in the bad-pair graph contains a triple of family SUM or of family SUM$'$.
\end{conjecture}

Conjecture~\ref{conj:link} implies the window statement for every $m$, since the prime case is
Theorem~\ref{thm:prime} and the two families are verified in
Proposition~\ref{prop:triples}. It has been checked for $4 \le m \le 30$ by a search that looks for
a counterexample directly rather than enumerating candidates, every run terminating without a cap.

Part of the conjecture is governed by classical extremal theory, and we record that because it
localises what is genuinely new. Nothing below is proved for a general $m$: the extremal input is
asymptotic with an unspecified threshold, so every statement in this passage is computed over the
range stated. A set is \emph{sum-free} if it contains no solution of $a+b=c$ with
$a,b,c$ in the set, $a=b$ permitted. The largest sum-free subsets of $[1,2m]$ have exactly $m$
elements, so a candidate $W$ has precisely the extremal size, and the extremal sets are classified.

\begin{proposition}[Computed for $5\le m\le12$; classical in substance]\label{prop:sf}
For $5\le m\le12$ the sum-free subsets of $[1,2m]$ of size $m$ are exactly three: the odd numbers,
the interval $[m,2m-1]$, and the interval $[m+1,2m]$. Each of them contains $1$ or $m+1$.
\end{proposition}

\begin{corollary}[Computed for $5\le m\le20$]\label{cor:schur}
For $5\le m\le20$, every $m$-element $W\subseteq[1,2m]$ avoiding $1$ and $m+1$ contains a solution of
$a+b=c$ with $a,b,c\in W$, where $a=b$ is permitted.
\end{corollary}

\begin{proof}
The odd numbers contain $1$, and both intervals contain $m+1$, so no set of the extremal size avoids
both; hence $W$ is not sum-free. The range $5\le m\le20$ has been verified directly, by searching for
a sum-free $m$-subset of $[1,2m]$ avoiding $1$ and $m+1$ and finding none. At $m=4$ the statement
fails: $\{2,3,7,8\}$ is sum-free, has four elements in $[1,8]$, and avoids $1$ and $5$.
\end{proof}

The case $a=b$ is not excluded. One might expect independence to rule it out, since $a=b$ puts
$\{a,2a\}$ inside $W$, but a doubling pair need not be a bad pair: MUL requires $2m/3<a\le m$ and
DBL requires $2a<m+3$, so the band $(m+3)/2<a\le2m/3$ is covered by neither, and there $\{a,2a\}$ is
an edge of no graph. The smallest instances are $a=9,10$ at $m=15$, $a=10,11$ at $m=17$ and
$a=12,13$ at $m=20$. So Corollary~\ref{cor:schur} yields a Schur solution but not necessarily three
distinct elements.

What Conjecture~\ref{conj:link} asks beyond Corollary~\ref{cor:schur} is that some sum triple have
distinct entries and meet the threshold $2c\ge3m+1$ and the coprimality conditions of the two
families. Existence of a solution is settled; distinctness and the side conditions are not. That is
where the remaining content sits.

We do not claim the extremal input as new. The maximum size $\lceil n/2\rceil$ of a sum-free subset
of $[1,n]$ is classical, and the structure of sum-free subsets of an interval near that size has been
studied at length; Tran~\cite{Tran} characterises the sum-free subsets of $[1,n]$ of density at
least $2/5-c$ for an absolute $c>0$, and our candidates have density exactly $1/2$. The alternative
in that characterisation that a large sum-free set is either entirely odd or has $\min A\ge|A|$
gives Proposition~\ref{prop:sf} for all large $m$, the two intervals being the two sum-free
$m$-subsets of $[m,2m]$. We verify the range we use rather than
track the exact threshold at which the structure theorem takes hold.

Four measurements delimit what a proof must do.

First, family DBL$_3$ is not needed: no independent set is met by DBL$_3$ alone at any $m$ checked,
while SUM$'$ alone meets $1$, $2$ and $9$ of them at $m = 10, 16, 22$. Both sum families are
required and the third is not.

Second, writing $c = a+b$, family SUM requires $a > 2m - c$ and SUM$'$ requires $a \le c-m-1$, so
between them they cover every splitting of $c$ exactly when $2c \ge 3m+1$, and below that threshold
the splitting $a = c-m$ is reached by neither. The threshold is therefore sharp.

Third, every candidate does contain sum triples above that threshold, between $3$ and $9$ of them,
with none lacking one for $4 \le m \le 30$. This is not a statement about sum-free sets alone: with
independence dropped, the block $\{2,3,\dots,m,m+2\}$ has all its sums too small and is a
counterexample at every $m$. That block carries every pair $\{a,2a\}$ with $2a \le m$, which is what
independence forbids.

Fourth, the coprimality conditions are severe. Of a candidate's sum triples above the threshold, as
few as $2$ of $36$ satisfy the coprimality that SUM or SUM$'$ also demands, at $m = 22$. The count
of usable triples never falls below $2$ for $m \ge 13$, but the proportion does not stay bounded
away from zero, so a counting argument over the threshold alone will not establish
Conjecture~\ref{conj:link}.

\section{Open problems}\label{sec:open}

\begin{enumerate}
\item The uniform window statement: does \eqref{eq:lrc} hold for every $m$ and every velocity set
inside $[1,2m]$? Equivalently, by Theorem~\ref{thm:equiv}, is $\tau(2m) = m+1$ for every $m$?
\item Classify the bad sets of size four and above. Sizes one to three are treated in
Section~\ref{sec:vocab}; the pair classification is now complete
(Theorem~\ref{thm:identity} and Corollary~\ref{cor:pairupper-uncond}), and the triples are handled
by families that Proposition~\ref{prop:triples-incomplete} shows are not exhaustive. Nothing is
known about size four. The empirical observation of item~4 below suggests they may not matter for
the decision procedure, which would be worth proving.
\item Classify the minimal bad triples that Proposition~\ref{prop:triples-incomplete} shows lie
outside SUM, SUM$'$ and DBL$_3$. The two examples given fail for different reasons, one on a
boundary of the hypotheses and one on a coprimality condition, which suggests the families need
widening rather than supplementing.
\item Is the empirical observation that irredundant cores use no bad set of size four or more true
for every $m$? If so, the window statement depends only on the sets classified in
Section~\ref{sec:vocab}.
\item Give an arithmetic description of the Gallai--Edmonds set $D$ of the bad-pair graph. All
candidates known to us fail, as recorded in Section~\ref{sec:obstructions}.
\end{enumerate}

\begin{thebibliography}{99}

\bibitem{Bed} B.~Bedert, \emph{Riesz products and the lonely runner conjecture: a wider gap of
loneliness}, 2025. arXiv:2511.16636.

\bibitem{Tran} T.~Tran, \emph{On the structure of large sum-free sets of integers}, Israel J.
Math. \textbf{228} (2018), 249--292. arXiv:1705.02584.

\bibitem{BHV} V.~Beresnevich, A.~Haynes, S.~Velani, \emph{Sums of reciprocals of fractional parts
and multiplicative Diophantine approximation}, Mem. Amer. Math. Soc. \textbf{263} (2020), no.~1276.

\bibitem{BP} T.~Bohman, F.~Peng, \emph{Coprime mappings and lonely runners}, Mathematika
\textbf{68} (2022), 784--804. arXiv:2109.09860.

\bibitem{Dem} R.~W.~Deming, \emph{Independence numbers of graphs: an extension of the
K\"onig--Egerv\'ary theorem}, Discrete Math. \textbf{27} (1979), 23--33.

\bibitem{Lar} C.~E.~Larson, \emph{The critical independence number and an independence
decomposition}, European J. Combin. \textbf{32} (2011), 294--300.

\bibitem{Pan} R.~K.~Pandey, \emph{On the lonely runner conjecture}, Math. Bohem. \textbf{135}
(2010), 63--68.

\bibitem{Ste} F.~Sterboul, \emph{A characterization of the graphs in which the transversal number
equals the matching number}, J. Combin. Theory Ser. B \textbf{27} (1979), 228--229.

\bibitem{Cassels} J.~W.~S.~Cassels, \emph{An Introduction to the Geometry of Numbers}, Springer,
1959.

\bibitem{CG} S.~Czerwi\'nski, J.~Grytczuk, \emph{Invisible runners in finite fields},
Inform. Process. Lett. \textbf{108} (2008), no.~2, 64--67.

\bibitem{Har} N.~M.~Haralambis, \emph{Sets of integers with missing differences},
J. Combin. Theory Ser. A \textbf{23} (1977), no.~1, 22--33.

\bibitem{MSS} R.~D.~Malikiosis, F.~Santos, M.~Schymura, \emph{Linearly-exponential checking is
enough for the Lonely Runner Conjecture and some of its variants}, Forum Math. Sigma \textbf{13}
(2025), e164, 32~pp. arXiv:2411.06903.

\bibitem{PS} G.~Perarnau, O.~Serra, \emph{The Lonely Runner Conjecture turns 60},
Comput. Sci. Rev. \textbf{58} (2025), 100798. arXiv:2409.20160.

\bibitem{Pom} C.~Pomerance, \emph{Coprime matchings}, INTEGERS \textbf{22} (2022), \#A2.

\bibitem{Ros8} M.~Rosenfeld, \emph{The lonely runner conjecture holds for eight runners}, 2025.
arXiv:2509.14111.

\bibitem{Ros9} M.~Rosenfeld, \emph{The lonely runner conjecture holds for nine runners}, 2025.
arXiv:2512.01912.

\bibitem{ST} T.~Sungkawichai, T.~Trakulthongchai, \emph{Eleven, twelve, and thirteen lonely
runners}, 2026. arXiv:2604.23906.

\bibitem{Tao} T.~Tao, \emph{Some remarks on the lonely runner conjecture}, Contrib. Discrete Math.
\textbf{13} (2018), no.~2, 1--31. arXiv:1701.02048.

\bibitem{Tra} T.~Trakulthongchai, \emph{Nine and ten lonely runners}, Electron. J. Combin.
\textbf{33} (2026), no.~2, \#P2.46. arXiv:2511.22427.

\end{thebibliography}

\end{document}
